\documentclass[10pt,a4paper]{amsart}
\usepackage[margin=19mm]{geometry}
\usepackage{amsmath,amssymb,mathtools}
\usepackage[scr=rsfs,cal=euler]{mathalfa}
\usepackage{xfrac}
\usepackage[unicode,hidelinks,bookmarks=false]{hyperref}
\newcommand{\ie}{i.e.\ }
\newcommand{\cf}{cf.\ }
\newcommand{\resp}{respectively\ }
\newcommand{\Subsection}{\subsection}
\providecommand{\nobreakdash}{\nobreak\hbox{-}}
\setlength{\emergencystretch}{2em}
\multlinegap0pt
\usepackage{bm}
\usepackage{bbm}
\usepackage{dsfont}
\usepackage{xspace}
\usepackage{cancel}
\usepackage{mathtools}
\usepackage{amsthm}
\newtheorem{theorem}{Theorem}
\newtheorem{lemma}[theorem]{Lemma}
\title[The Marcus-Minc Transform Inequality]{The Marcus-Minc Transform Inequality (Lemma 3.1)}
\author{Yair Lavi}
\date{Source: arXiv:2609.29262v1 (CC BY 4.0)}
\begin{document}
\maketitle

\noindent\emph{Source and licence.} The Marcus-Minc Transform Inequality. Version arXiv:2609.29262v1 (CC BY 4.0), distributed under the Creative Commons Attribution 4.0 International licence, as recorded in the arXiv licence metadata for this exact version and shown on the abstract page https://arxiv.org/abs/2609.29262v1. Full rights notice: licenses/arxiv-2609.29262-CC-BY-4.0.txt.\par\smallskip

\noindent\textit{Editorial scope.} Counted unit: the Lemma printed as Lemma 3.1 of the article (source0.tex lines 401--410, source label \texttt{None}), delivered together with the article's own complete original proof of it (source0.tex lines 412--537, one \texttt{proof} environment of 126 lines). No mathematics is written, repaired or re-derived: statement and proof are the article's own text line for line. Required context retained uncounted: none. Every definition, display and label of those lines is retained unchanged. Omitted from this package and not redistributed: 1--400, 411--411, 538--824. Nothing inside a retained excerpt is omitted or altered: every retained body is delivered exactly as the article's lines are written, so no omission receipt of any kind accompanies this package. Citations: the retained excerpts carry no citation command at all, so no bibliography entry of the article is transcribed and no reference is redistributed. Reference adaptation: none was needed and none was made; every reference inside the delivered unit resolves inside it, so no unresolved reference or citation is produced. Printed slips kept literal: no printed word, symbol, hypothesis or number of the counted statement or of its proof was altered. Layout and class: the article's own class is \texttt{article} with packages including fontenc, lmodern, amsmath, amssymb, amsthm, geometry, cite, hyperref; the wrapper is the lane's pinned \texttt{amsart} editorial wrapper at 10pt on A4, which restates the theorem-like environments the retained text needs and adds the article's own macro definitions that the retained text needs (none), with extra wrapper packages (bm, bbm, dsfont, xspace, cancel, mathtools). The delivered numbering is the unit's own. Assets: no retained span contains a figure environment, a graphics-inclusion command, a PDF-page inclusion, a table or a TikZ picture; no image file is copied into this package, the package declares no asset dependency and no image file is redistributed with it. Licence: version arXiv:2609.29262v1 of this article is distributed under the Creative Commons Attribution 4.0 International licence, as recorded in the arXiv licence metadata for this exact version and shown on the abstract page https://arxiv.org/abs/2609.29262v1. A full rights notice is delivered at licenses/arxiv-2609.29262-CC-BY-4.0.txt. Characters: every retained excerpt is pure ASCII, so the delivered engine, XeLaTeX with its default Latin Modern text font, reports no missing character.
\begin{lemma}
If $X$ is a real $n\times n$ matrix with zero margins and
$\|X\|_{2\to2}\le1$, then, for
$2\le k\le n$,
\[
|\sigma_k(X)|\le
\frac{\sigma_k(Q)}{n-1}\|X\|_F^2.
\tag{3.2}
\]
\end{lemma}

\begin{proof}
We use tensor powers of $\mathbb R^n$ with the product inner
product
\[
\langle x_1\otimes\cdots\otimes x_k,
y_1\otimes\cdots\otimes y_k\rangle
=\prod_{\ell=1}^k\langle x_\ell,y_\ell\rangle,
\]
extended bilinearly. If $S=\{i_1,\ldots,i_k\}$ is a $k$-element subset of
$[n]$, put
\[
v_S=\frac1{\sqrt{k!}}\sum_{\pi\in S_k}
e_{i_{\pi(1)}}\otimes\cdots\otimes e_{i_{\pi(k)}},
\qquad
u_k=\sum_{|S|=k}v_S,
\tag{3.3}
\]
where $e_1,\ldots,e_n$ are the standard basis vectors. The vectors $v_S$
are orthonormal, and expanding the inner product gives
\[
\langle v_S,X^{\otimes k}v_T\rangle=\operatorname{per}X[S,T].
\tag{3.4}
\]
Zero margins imply $X=QXQ$. Hence, with
$w=Q^{\otimes k}u_k\in H^{\otimes k}$,
\[
\sigma_k(X)=\langle w,X^{\otimes k}w\rangle,
\qquad
\|w\|^2=\sigma_k(Q)\ge0.
\tag{3.5}
\]
The second identity follows because $Q^{\otimes k}$ is an orthogonal
projection. It also supplies the nonnegativity of every $\sigma_k(Q)$ used
below.

We now define the tensor property responsible for the estimate. A linear
operator $R$ on $H$ is \textbf{self-adjoint} if
$\langle Rx,y\rangle=\langle x,Ry\rangle$ for all $x,y\in H$. For
$w\in H^{\otimes k}$ and $1\le\ell\le k$, its $\ell$-th \textbf{one-factor
marginal} is the unique self-adjoint operator $\rho_\ell$ on $H$ satisfying
\[
\operatorname{tr}(\rho_\ell R)=
\left\langle w,
\bigl(I_H^{\otimes(\ell-1)}\otimes R\otimes
I_H^{\otimes(k-\ell)}\bigr)w\right\rangle
\tag{3.6}
\]
for every self-adjoint operator $R$ on $H$; here $\operatorname{tr}$ denotes
the sum of the diagonal entries of an operator in an orthonormal basis. We
call $w$ \textbf{one-factor isotropic} when each $\rho_\ell$ is a scalar multiple
of $I_H$.

To see existence, choose an orthonormal basis $(h_a)$ of $H$ and, after putting
the $\ell$-th factor first, write $w=\sum_a h_a\otimes z_a$, where
$z_a\in H^{\otimes(k-1)}$. The matrix
$(\langle z_a,z_b\rangle)_{a,b}$ is self-adjoint and satisfies (3.6), so it
defines $\rho_\ell$. If two self-adjoint operators satisfy (3.6), their
difference $D$ gives $\operatorname{tr}(D^2)=0$ upon taking $R=D$, and therefore
$D=0$.

Simultaneously permuting the $n$ coordinate labels fixes $u_k$, commutes with
$Q$, and therefore fixes $w$. Thus each $\rho_\ell$ commutes with every
coordinate permutation on $H$. Extend $\rho_\ell$ by zero on
$\mathbb R\mathbf1$. A matrix commuting with every permutation has one
constant on its diagonal and another off its diagonal. Its restriction to
$H$ is therefore scalar. Since
$\operatorname{tr}\rho_\ell=\|w\|^2$, we obtain
\[
\rho_\ell=\frac{\|w\|^2}{n-1}I_H
\qquad(1\le\ell\le k).
\tag{3.7}
\]
This is the asserted isotropy; it is a property of $w$, not an assumption on
$X$.

Choose linear operators $U,V$ on $H$ satisfying
$U^{\mathsf T}U=V^{\mathsf T}V=I_H$ and numbers
$0\le s_1,\ldots,s_{n-1}\le1$ such that
\[
X|_H=UDV^{\mathsf T},
\qquad
D=\operatorname{diag}(s_1,\ldots,s_{n-1}).
\]
Such a factorization is a singular-value decomposition; the displayed
identities say that $U$ and $V$ are orthogonal, and the upper bounds on the
$s_i$ follow from $\|X\|_{2\to2}\le1$. Put
\[
a=(U^{\mathsf T})^{\otimes k}w,
\qquad
b=(V^{\mathsf T})^{\otimes k}w,
\qquad
K=D^{\otimes k}.
\]
The one-factor marginals of both $a$ and $b$ still equal (3.7). A self-adjoint
operator $R$ is \textbf{positive semidefinite} if
$\langle z,Rz\rangle\ge0$ for every $z$. For self-adjoint operators $R,S$,
write $R\preceq S$ when $S-R$ is positive semidefinite. For every diagonal
entry of $K$, the assumptions $k\ge2$ and $0\le s_i\le1$ give
\[
\prod_{\ell=1}^k s_{i_\ell}
\le\left(\prod_{\ell=1}^k s_{i_\ell}^2\right)^{1/k}
\le\frac1k\sum_{\ell=1}^k s_{i_\ell}^2.
\]
Equivalently,
\[
0\preceq K\preceq\frac1k\sum_{\ell=1}^k
I_H^{\otimes(\ell-1)}\otimes D^2\otimes
I_H^{\otimes(k-\ell)}.
\tag{3.8}
\]
Equations (3.6)--(3.8) yield
\[
\langle a,Ka\rangle,\ \langle b,Kb\rangle
\le\frac{\|w\|^2}{n-1}\operatorname{tr}D^2.
\]
Cauchy--Schwarz for the positive semidefinite operator $K$ now gives
\[
\begin{aligned}
|\sigma_k(X)|
&=|\langle a,Kb\rangle|\\
&\le\sqrt{\langle a,Ka\rangle\langle b,Kb\rangle}\\
&\le\frac{\sigma_k(Q)}{n-1}\|X\|_F^2,
\end{aligned}
\]
because $\operatorname{tr}D^2=\|X\|_F^2$. This proves (3.2).
\end{proof}

\end{document}
